% o README
%
%     These  rulers have  been  designed using  LaTeX  and the  PGFPlots
%     package by  Christian Feuersanger to  measure slope on  a 1:25,000
%     scale or 1:50,000 topographic map.
%
%     Slope  is determined  by  comparing the  distance between  contour
%     lines of  the map  with the  reference lines on  the edges  of the
%     ruler.  The top  border of the ruler is  designed to measure slope
%     with  a 20~meter  contour  interval.  The  bottom border  measures
%     slope with a 50~meter contour interval.
%
%     You should print this ruler  on white or transparent paper, taking
%     care of  deactivating any scaling options.  Check  the actual size
%     of the ruler in centimeters on the right border.
%
% o COPYRIGHT NOTICE
%
%     Copyright (C) 2013 Emmanuel Vazquez
%
%     Version: 2.0
%
%     Author:  Emmanuel Vazquez  <emmanuel.vazquez@free.fr>
%
%     URL:
%
% o COPYING PERMISSION STATEMENT
% 
%    This file is licensed under a Creative Commons Attribution -
%    NonCommercial 3.0 Unported License.
%
%    * You are free:
%      - to Share      -- to copy, distribute and transmit the work
%      - to Remix      -- to adapt the work
%
%    * Under the following conditions:
%      - Attribution  --  You must  attribute  the  work  in the  manner
%                         specified by  the author or  licensor (but not
%                         in any way that suggests that they endorse you
%                         or your use of the work).
%      - Noncommercial -- You may not use this work for commercial
%                         purposes.
%%%%%%%%%%%%%%%%%%%%%%%%%%%%%%%%%%%%%%%%%%%%%%%%%%%%%%%%%%%%%%%%%%%%%%%%

\documentclass[a4paper,10pt]{article}

\usepackage[latin1]{inputenc}
\usepackage{arev}
\usepackage[T1]{fontenc}
\usepackage{pgfplots}
\usepackage{amsmath,amsfonts}
\usepackage{bm, bbm}

\begin{document}
\pagestyle{empty}

%%%%%%%%%%%%%%%%%%%%%%%%%%%%%%%%%%%%%%%%%%%%%%%%%%%%%%%%%%%%%%%%%%%%%%%%

\def\printf#1{\pgfkeys{/pgf/number format/.cd,fixed,precision=2,zerofill=true}\pgfmathparse{#1}\pgfmathprintnumber{\pgfmathresult}}
\def\printi#1{\pgfkeys{/pgf/number format/.cd,fixed,precision=0}\pgfmathparse{#1}\pgfmathprintnumber{\pgfmathresult}}

\def\framew{8.5cm}
\def\frameh{5.4cm}
\def\lineh {1.2cm}
\def\sline {+(0,0) -- +(0,\lineh)}

%%%%%%%%%%%%%%%%%%%%%%% 1:25,000 slope ruler %%%%%%%%%%%%%%%%%%%%%%%%%%%

\def\scale{25000}
\def\dl{20}
\def\dh{50}
\def\spacings {{0cm, 1.2cm, 2.4cm, 3.5cm, 4.6cm}}

\begin{center}
\begin{tikzpicture}
  % draw the frame
  \draw[gray, very thin] (0,0) rectangle (\framew,\frameh);

  % l=20m ruler
  \pgfkeys{/pgf/fpu=true,/pgf/fpu/output format=fixed}
  \foreach \i/\th in {0/45, 1.1/40, 2.2/35, 3.3/30, 4.5/25}
     \draw[very thin] (15mm+\i cm, \frameh - \lineh) \sline ++({\dl/tan(\th)/\scale*100},0) \sline
                     +(0, 0.2cm) node[anchor=west] {\scriptsize \printi{\th}$^{\circ}$}
                     +(0, 0.5cm) node[anchor=west] {\scriptsize \printf{\dl/tan(\th)/\scale*1000}};

  % l=50m ruler
  \foreach \i/\th in {0/45, 1.1/40, 2.2/35, 3.3/30, 4.5/25}
     \draw[very thin] (15mm+\i cm, 0cm) \sline ++({\dh/tan(\th)/\scale*100},0) \sline
                     +(0, \lineh - 0.2cm) node[anchor=west] {\scriptsize \printi{\th}$^{\circ}$}
                     +(0, \lineh - 0.5cm) node[anchor=west] {\scriptsize \printf{\dh/tan(\th)/\scale*1000}};
  \pgfkeys{/pgf/fpu=false}

  % right border: cm scale
  \def\x{(\frameh-4cm)/2}
  \foreach \dx in {0,0.1,...,4}
    \draw (\framew, {\x + \dx cm}) -- +(-1mm, 0); 
  \foreach \dx in {0,0.5,...,4}
    \draw (\framew, {\x + \dx cm}) -- +(-2mm, 0);
  \foreach \dx in {0,...,4}
    \draw ({\framew-2mm}, {\x + \dx cm}) node[anchor=east] {$\dx$};
  \draw (\framew,\frameh-3mm) node[anchor=east]{\scriptsize cm};

  % left border: km scale
  \pgfkeys{/pgf/fpu=true,/pgf/fpu/output format=fixed}
  \foreach \dx in {0,0.05,...,1}
    \draw[very thin] (0, {\x + \dx cm/\scale*1e5}) -- +(0.7mm, 0); 
  \foreach \dx in {0,0.1,...,1}
    \draw (0, {\x + \dx cm/\scale*1e5}) -- +(2mm, 0); 
  \foreach \dx in {0,0.5,...,1}
    \draw (0, {\x + \dx cm/\scale*1e5}) -- +(3mm, 0);
  \foreach \dx in {0,0.5,...,1}
    \draw (3mm, {\x + \dx cm/\scale*1e5}) node[anchor=west] {$\dx$};
  \draw (0,\frameh-3mm) node[anchor=west]{\scriptsize km};
  \pgfkeys{/pgf/fpu=false}

  % information about the ruler
  \draw (13.5mm,\frameh/2) node[anchor=west] {\scriptsize $1:\scale$ Scale
    Slope Ruler};
  \draw (15mm, \frameh/2 + 12mm) -- ++(0, -2mm) -- 
    ++(57mm,0) -- ++(0, +2mm) ;
  \draw (\framew/2, \frameh/2 +8mm) node {\scriptsize $h=\dl\rm m$};

  \draw (15mm, \frameh/2 - 12mm) -- ++(0, +2mm) -- 
    ++(57mm,0) -- ++(0, -2mm) ;
  \draw (\framew/2, \frameh/2 -8mm) node {\scriptsize $h=\dh\rm m$};
  \pgftext[at=\pgfpoint{5.7cm}{\frameh/2+4mm},left] {\includegraphics[width=1.5cm]{logoc2c}};
  \pgftext[at=\pgfpoint{5.7cm}{\frameh/2-4mm},left] {\includegraphics[width=1.5cm]{by-nc}};
  % 0.8mm control lines
  \draw[very thin] (15mm, 4.0cm) -- +(0, 2mm) ++(0.8mm, 0) -- +(0, 2mm);

\end{tikzpicture}

\vskip 1cm

%%%%%%%%%%%%%%%%%%%%%%% 1:50,000 slope ruler %%%%%%%%%%%%%%%%%%%%%%%%%%%

\def\scale{50000}
\def\dl{20}
\def\dh{50}
\def\spacings {{0cm, 1.2cm, 2.4cm, 3.5cm, 4.6cm}}

\noindent
\begin{tikzpicture}
  % draw the frame
  \draw[gray, very thin] (0,0) rectangle (\framew,\frameh);

  % l=20m ruler
  \pgfkeys{/pgf/fpu=true,/pgf/fpu/output format=fixed}
  \foreach \i/\th in {0/45, 1.1/40, 2.2/35, 3.3/30, 4.5/25}
     \draw[very thin] (15mm+\i cm, \frameh - \lineh) \sline ++({\dl/tan(\th)/\scale*100},0) \sline
                     +(0, 0.2cm) node[anchor=west] {\scriptsize \printi{\th}$^{\circ}$}
                     +(0, 0.5cm) node[anchor=west] {\scriptsize \printf{\dl/tan(\th)/\scale*1000}};

  % l=50m ruler
  \foreach \i/\th in {0/45, 1.1/40, 2.2/35, 3.3/30, 4.5/25}
     \draw[very thin] (15mm+\i cm, 0cm) \sline ++({\dh/tan(\th)/\scale*100},0) \sline
                     +(0, \lineh - 0.2cm) node[anchor=west] {\scriptsize \printi{\th}$^{\circ}$}
                     +(0, \lineh - 0.5cm) node[anchor=west] {\scriptsize \printf{\dh/tan(\th)/\scale*1000}};
  \pgfkeys{/pgf/fpu=false}

  % right border: cm scale
  \def\x{(\frameh-4cm)/2}
  \foreach \dx in {0,0.1,...,4}
    \draw (\framew, {\x + \dx cm}) -- +(-1mm, 0); 
  \foreach \dx in {0,0.5,...,4}
    \draw (\framew, {\x + \dx cm}) -- +(-2mm, 0);
  \foreach \dx in {0,...,4}
    \draw ({\framew-2mm}, {\x + \dx cm}) node[anchor=east] {$\dx$};
  \draw (\framew,\frameh-3mm) node[anchor=east]{\scriptsize cm};

  % left border: km scale
  \pgfkeys{/pgf/fpu=true,/pgf/fpu/output format=fixed}
  \foreach \dx in {0,0.05,...,2}
    \draw[very thin] (0, {\x + \dx cm/\scale*1e5}) -- +(0.7mm, 0); 
  \foreach \dx in {0,0.1,...,2}
    \draw (0, {\x + \dx cm/\scale*1e5}) -- +(2mm, 0); 
  \foreach \dx in {0,0.5,...,2}
    \draw (0, {\x + \dx cm/\scale*1e5}) -- +(3mm, 0);
  \foreach \dx in {0,0.5,...,2}
    \draw (3mm, {\x + \dx cm/\scale*1e5}) node[anchor=west] {$\dx$};
  \draw (0,\frameh-3mm) node[anchor=west]{\scriptsize km};
  \pgfkeys{/pgf/fpu=false}

  % information about the ruler
  \draw (13.5mm,\frameh/2) node[anchor=west] {\scriptsize $1:\scale$ Scale
    Slope Ruler};
  \draw (15mm, \frameh/2 + 12mm) -- ++(0, -2mm) -- 
    ++(57mm,0) -- ++(0, +2mm) ;
  \draw (\framew/2, \frameh/2 +8mm) node {\scriptsize $h=\dl\rm m$};

  \draw (15mm, \frameh/2 - 12mm) -- ++(0, +2mm) -- 
    ++(57mm,0) -- ++(0, -2mm) ;
  \draw (\framew/2, \frameh/2 -8mm) node {\scriptsize $h=\dh\rm m$};
  \pgftext[at=\pgfpoint{5.7cm}{\frameh/2+4mm},left] {\includegraphics[width=1.5cm]{logoc2c}};
  \pgftext[at=\pgfpoint{5.7cm}{\frameh/2-4mm},left] {\includegraphics[width=1.5cm]{by-nc}};
  % 0.4mm control lines
  \draw[very thin] (15mm, 4.0cm) -- +(0, 2mm) ++(0.4mm, 0) -- +(0, 2mm);

\end{tikzpicture}
\end{center}


%%%%%%%%%%%%%%%%%%%%%%%%%%%%%%%%%%%%%%%%%%%%%%%%%%%%%%%%%%%%%%%%%%%%%%%% 
\vskip 1.0cm
\small

\noindent
These rulers have been designed  using \LaTeX{} and the PGFPlots package
by Christian Feuers\"anger to measure  slope on a 1:25,000 or a 1:50,000
scale  topographic map. Slope  is determined  by comparing  the distance
between contour lines  of the map with the reference  lines on the edges
of the ruler.  The top border  of the ruler is designed to measure slope
with a 20~meter contour interval.  The bottom border measures slope with
a 50~meter contour interval.  \bigskip

\noindent
You should print  this ruler on white or  transparent paper, taking care
of deactivating any scaling options.  Check the actual size of the ruler
in centimeters on the right border. \bigskip

\noindent   \sl  These   C2C   Slope  Rulers   are   licensed  under   a
Creative~Commons Attribution  -- NonCommercial 3.0  Unported License (CC
BY-NC 3.0) \bigskip

\noindent
v2.0 Emmanuel Vazquez
\end{document}

%%% Local Variables:
%%% mode:latex
%%% End:
